% ─────────────────────────────────────────────────────────────────────────────
%
% Documentación de la estimación del plano del suelo por RANSAC jerárquico.
% Corresponde a la tercera iteración del módulo local_mapper.
% ─────────────────────────────────────────────────────────────────────────────

\section{Detección del plano del suelo}
\label{sec:lm-ground-estimator}

Tener una nube de puntos con el eje $Y$ alineado a la dirección de la
gravedad, es una corrección necesaria para poder interpretar la
geometría de la escena. El suelo, si es plano, debe aparecer ahora como un
conjunto de puntos coplanares, todos a la misma altura $Y \approx h$, donde
$h$ es la distancia de la cámara al suelo. Esta iteración introduce la
detección y parametrización de ese plano del suelo mediante un pipeline de
ocho etapas que sigue la metodología propuesta por Poux et~al.\ \cite{Poux2019voxel, Poux2020segmentation} para el
procesamiento robusto de nubes de puntos 3D.

\subsection{Marco teórico}

\subsubsection{Representación de un plano en 3D}

Un plano en el espacio tridimensional se puede representar mediante su
ecuación implícita:
\begin{equation}
  \label{eq:plane-implicit}
  \mathbf{n} \cdot \mathbf{p} + d = 0
\end{equation}
donde $\mathbf{n} = (n_x, n_y, n_z)^T \in \mathbb{R}^3$ es el vector
normal unitario al plano ($\lVert\mathbf{n}\rVert = 1$),
$\mathbf{p} = (x, y, z)^T$ es un punto arbitrario del plano, y $d \in
\mathbb{R}$ es la distancia perpendicular con signo del origen de coordenadas al
plano. La distancia con signo de un punto arbitrario $\mathbf{q}$ al plano es:
\begin{equation}
  \label{eq:point-plane-dist}
  \delta(\mathbf{q}) = \mathbf{n} \cdot \mathbf{q} + d
\end{equation}
Esta expresión es positiva si $\mathbf{q}$ está del mismo lado que
$\mathbf{n}$, y negativa en caso contrario. El origen de coordenadas
(posición de la cámara) se encuentra a distancia $|d|$ del plano.

\subsubsection{Estimación de ruido por distancia al vecino más cercano}
\label{subsec:lm-v3-noise}

El nivel de ruido de la cámara de profundidad varía con la distancia y con
la reflectividad de la superficie. Para adaptar automáticamente la tolerancia
de inliers sin intervención manual, se calcula estadísticamente la distancia
media al vecino más cercano ($\bar{d}_{NN}$) sobre un subconjunto de la
nube \cite{gonzalez2018digital}:
\begin{equation}
  \label{eq:nn-dist}
  \bar{d}_{NN} = \frac{1}{m} \sum_{i=1}^m \min_{j \neq i} \lVert \mathbf{p}_i - \mathbf{p}_j \rVert
\end{equation}
La búsqueda del vecino más cercano para cada uno de los $m$ puntos muestreados
se realiza como una búsqueda exhaustiva sobre la propia muestra, con
complejidad $O(m^2)$. Si bien existen estructuras de datos más eficientes
(kd-tree, $O(m \log m)$), la búsqueda exhaustiva es adecuada para el tamaño
de muestra utilizado ($m = 200$ por defecto, parámetro
\texttt{nn\_sample\_size}): el costo absoluto es de $200^2 = 40\,000$
comparaciones de distancia, con latencia despreciable respecto al conjunto
del pipeline. La muestra se selecciona mediante una permutación aleatoria
con semilla fija, lo que garantiza reproducibilidad entre frames consecutivos
y facilita el análisis y el testing: dos nubes con el mismo contenido producen
siempre las mismas estadísticas de diagnóstico.

La desviación estándar de esta distribución, $\sigma_{NN}$, sirve como
estimación del ruido puntual del sensor. Bajo la hipótesis de que el ruido
de medición sigue una distribución aproximadamente gaussiana, el intervalo
$\pm 3\,\sigma_{NN}$ contiene el $99{,}7\,$\% de los valores esperados para
puntos que genuinamente pertenecen al plano. Fijar la tolerancia en ese valor
garantiza que casi ningún punto válido sea descartado como outlier por error
de medición, a la vez que rechaza perturbaciones significativamente mayores.
La tolerancia adaptativa se establece entonces como:

\begin{equation}
  \label{eq:adaptive-tol}
  \varepsilon = 3 \, \sigma_{NN}
\end{equation}

acotada entre $1\,\text{cm}$ y $10\,\text{cm}$ para garantizar
estabilidad numérica ante nubes muy densas o muy dispersas.

\subsubsection{Downsampling por voxel grid}
\label{subsec:lm-v3-voxel}

Una nube de puntos generada por la cámara D435 a plena resolución contiene
del orden de $4 \times 10^5$ puntos por fotograma ($848 \times 480$ píxeles).
El tiempo de ejecución de RANSAC crece linealmente con el número de puntos
(etapa de conteo de inliers), por lo que es conveniente reducir la densidad
antes de ejecutar el algoritmo.

El \emph{voxel grid} es una técnica estándar de submuestreo espacial
\cite{hornung2013octomap}: se divide el espacio en una rejilla de celdas
cúbicas de lado $s$ y se reemplaza cada celda que contenga al menos un punto
por el centroide de todos los puntos que caen en ella. El número de puntos de
salida es proporcional al volumen de la escena dividido por $s^3$, con
independencia de la densidad de muestreo original. Este proceso reduce la
distorsión causada por regiones sobredensas (por ejemplo, superficies
cercanas a la cámara) y hace que la estimación de planos sea más uniforme
espacialmente.

Con la configuración por defecto ($s = 0{,}05\,\text{m}$), en escenas típicas
de interior la nube pasa de $\approx 262\,000$ puntos con profundidad válida a
$\approx 17\,000$ tras el downsampling, una reducción de aproximadamente
$15\times$. Nótese que la entrada al voxel grid ($\approx 262\,000$) es menor
que la resolución nominal del sensor ($848 \times 480 = 407\,040$~píxeles)
porque el sensor D435 devuelve profundidad inválida en píxeles sin retorno
suficiente (superficies muy reflectivas, bordes de oclusión, zonas fuera del
rango operativo).

\subsubsection{RANSAC para detección robusta de planos}
\label{subsec:lm-v3-ransac}

Un ajuste directo por mínimos cuadrados solo es apropiado cuando todos los
puntos pertenecen al plano; en una nube 3D real el plano del suelo coexiste
con puntos de obstáculos, paredes y ruido. El algoritmo RANSAC (\emph{Random
Sample Consensus}) proporciona una estimación robusta de modelos paramétricos
en presencia de valores atípicos (\emph{outliers}) \cite{Thrun2005}.

El ciclo básico de RANSAC para planos es:
\begin{enumerate}
  \item Muestrear aleatoriamente el número mínimo de puntos necesarios para
        ajustar el modelo (tres puntos definen un plano).
  \item Ajustar el plano a esos tres puntos por producto vectorial.
  \item Contar los puntos de la nube completa cuya distancia al plano
        candidato es menor que una tolerancia $\varepsilon$ (inliers).
  \item Retener el plano con mayor número de inliers.
  \item Repetir durante $K$ iteraciones.
\end{enumerate}
Con $K$ iteraciones suficientes, la probabilidad de no haber muestreado al
menos una vez un conjunto de tres puntos todos inliers es exponencialmente
pequeña.

Si la fracción de inliers en la nube es $\rho$, la probabilidad de que una
muestra aleatoria de tres puntos esté formada íntegramente por inliers es
$\rho^3$, asumiendo muestreo uniforme e independencia aproximada entre
extracciones. En consecuencia, la probabilidad de que una iteración de
RANSAC falle en producir una hipótesis libre de outliers es $1 - \rho^3$.
Bajo la misma aproximación, la probabilidad de fallar en las $K$
iteraciones queda dada por la ecuación~\ref{eq:ransac-failure-prob}:
\begin{equation}
  \label{eq:ransac-failure-prob}
  P_{\mathrm{fail}} = (1 - \rho^3)^K
\end{equation}
Esta expresión corresponde a la formulación estándar de RANSAC para un
modelo definido por una muestra mínima de tamaño $s = 3$ \cite{Thrun2005}.
Para $\rho = 0{,}5$ y $K = 100$, la ecuación~\ref{eq:ransac-failure-prob}
da una probabilidad de fallo inferior a $10^{-9}$.

\subsubsection{Ajuste de plano por mínimos cuadrados}
\label{subsec:lm-v3-ls}

Dado un conjunto de $n$ puntos $\{\mathbf{p}_i\}_{i=1}^n$ que se suponen
coplanares con ruido, se busca el plano $(\mathbf{n}, d)$ con
$\lVert\mathbf{n}\rVert = 1$ que minimiza la suma de cuadrados de las
distancias ortogonales. Usando la ecuación~\ref{eq:point-plane-dist}, la
distancia del punto $\mathbf{p}_i$ al plano candidato es
$\delta_i = \mathbf{n}^T\mathbf{p}_i + d$, por lo que la función de costo
$J$ (el total del error cuadrático) queda:
\begin{equation}
  \label{eq:ls-objective}
  J(\mathbf{n}, d)
  = \sum_{i=1}^n \delta_i^2
  = \sum_{i=1}^n \bigl(\mathbf{n}^T \mathbf{p}_i + d\bigr)^2
\end{equation}
y el problema es encontrar $\mathbf{n}$ y $d$ que lo hagan mínimo:
\begin{equation}
  \label{eq:ls-min}
  \min_{\mathbf{n},\,d} \; J(\mathbf{n}, d), \qquad \lVert\mathbf{n}\rVert = 1
\end{equation}
Derivando respecto a $d$ e igualando a cero se obtiene el centroide:
\begin{equation}
  \label{eq:centroid}
  \bar{\mathbf{p}} = \frac{1}{n} \sum_{i=1}^n \mathbf{p}_i, \qquad
  d^* = -\mathbf{n}^T \bar{\mathbf{p}}
\end{equation}
Sustituyendo, el problema se reduce a minimizar únicamente sobre
$\mathbf{n}$:
\begin{equation}
  \label{eq:rayleigh-pre}
  J\bigl(\mathbf{n}, d^*(\mathbf{n})\bigr)
  = \sum_{i=1}^n \bigl(\mathbf{n}^T (\mathbf{p}_i - \bar{\mathbf{p}})\bigr)^2
  = \mathbf{n}^T S\, \mathbf{n}
\end{equation}
donde en la última igualdad se usó que, para cada término,
$(\mathbf{n}^T \mathbf{a})^2 = \mathbf{n}^T (\mathbf{a}\mathbf{a}^T)\mathbf{n}$,
con $\mathbf{a} = \mathbf{p}_i - \bar{\mathbf{p}}$, y por tanto
\begin{equation}
  \label{eq:scatter}
  S = \sum_{i=1}^n
      (\mathbf{p}_i - \bar{\mathbf{p}})
      (\mathbf{p}_i - \bar{\mathbf{p}})^T
\end{equation}
es la matriz de dispersión de los puntos centrados. La matriz de
covarianza $3 \times 3$ es:
\begin{equation}
  \label{eq:covariance}
  C = \frac{1}{n} \sum_{i=1}^n
      (\mathbf{p}_i - \bar{\mathbf{p}})
      (\mathbf{p}_i - \bar{\mathbf{p}})^T
\end{equation}
de modo que $S = nC$. Como $n > 0$ es una constante, minimizar
$\mathbf{n}^T S\,\mathbf{n}$ o minimizar $\mathbf{n}^T C\,\mathbf{n}$
produce exactamente el mismo vector normal óptimo.
$C$ es semidefinida positiva y simétrica, por lo que admite la
descomposición espectral $C = Q \Lambda Q^T$ con autovalores
$\lambda_1 \geq \lambda_2 \geq \lambda_3 \geq 0$. El cociente de Rayleigh
$\mathbf{n}^T C\, \mathbf{n}$ alcanza su mínimo $\lambda_3$ cuando
$\mathbf{n}$ es el autovector $\mathbf{q}_3$ correspondiente al autovalor
más pequeño \cite[App.~A4]{Hartley2004}. En consecuencia, la solución óptima es:
\begin{equation}
  \label{eq:plane-d}
  \mathbf{n}^* = \mathbf{q}_3, \qquad
  d^* = -\mathbf{n}^{*T} \bar{\mathbf{p}}
\end{equation}
Este autovector indica la dirección de \emph{menor} varianza de la nube, que
en un conjunto aproximadamente coplanar coincide con la perpendicular al
plano.

En la implementación, el cálculo de autovectores de $C$ se realiza mediante
el algoritmo iterativo de Jacobi para matrices simétricas $3 \times 3$
\cite{stroustrup2013cpp}. La implementación trabaja sobre la matriz
simétrica acumulada de segundo orden y realiza barridos sucesivos sobre los
tres pares de índices $(0,1)$, $(0,2)$ y $(1,2)$. Para cada par $(i,j)$,
si el término fuera de la diagonal es no nulo, se construye una rotación
ortogonal $R_{ij}(\theta)$ y se aplica la transformación de similitud
\begin{equation}
  \label{eq:jacobi-rotation}
  C_{k+1} = R_{ij}(\theta)^T\, C_k\, R_{ij}(\theta)
\end{equation}
con el objetivo de anular el elemento no diagonal $(C_k)_{ij}$. En la
implementación, el ángulo se calcula como
\begin{equation}
  \label{eq:jacobi-angle}
  \theta = \frac{1}{2}
  \operatorname{atan2}\!\left(2(C_k)_{ij}, (C_k)_{ii} - (C_k)_{jj}\right)
\end{equation}
de modo que la submatriz $2 \times 2$ asociada a las filas y columnas
$i,j$ quede diagonalizada tras la rotación. El proceso se repite durante un
número fijo de barridos; en cada uno se actualiza tanto la matriz simétrica
como la matriz acumulada de rotaciones, cuyas columnas convergen a los
autovectores. Finalmente, la normal del plano se toma como el autovector
asociado al menor autovalor. En este problema, al tratarse de matrices
$3 \times 3$, el costo computacional es bajo y no se requieren dependencias
externas.

\subsection{Arquitectura}

Se extiende el pipeline del nodo con la clase \texttt{GroundEstimator}.
Esta clase recibe como entrada la nube de puntos ya alineada a gravedad
(frame donde $Y$ apunta hacia abajo) y produce como salida el plano del
suelo estimado, junto con listas de índices clasificando cada punto como
suelo (\texttt{GROUND}), obstáculo (\texttt{OBSTACLE}) o techo
(\texttt{CEILING}). El resto del módulo (DepthProjector, ImuFilter, GravityAligner)
se mantiene sin cambios.


\subsection{Método: pipeline de ocho etapas}
\label{sec:lm-v3-pipeline}

El pipeline de \texttt{GroundEstimator} sigue ocho etapas secuenciales sobre
cada fotograma de profundidad.

\paragraph{Etapa 1: Diagnóstico de la nube.}
Se calcula sobre una muestra aleatoria de $m$ puntos (por defecto $m = 200$)
la distancia al vecino más cercano, su media $\bar{d}_{NN}$ y desviación
estándar $\sigma_{NN}$ (ecuación~\ref{eq:nn-dist}). Adicionalmente se estima
la densidad volumétrica de la nube y el porcentaje de puntos outliers
(aquellos cuya distancia NN supera $\bar{d}_{NN} + 3\,\sigma_{NN}$). Estas
estadísticas se exponen públicamente para monitoreo y ajuste de parámetros
en tiempo de ejecución.

\paragraph{Etapa 2: Submuestreo por voxel grid.}
La nube completa se reduce mediante un voxel grid de resolución $s =
0{,}05\,\text{m}$ por defecto (sección~\ref{subsec:lm-v3-voxel}). Se utiliza
una tabla hash con clave $(i_x, i_y, i_z) = \lfloor \mathbf{p}/s \rfloor$
para acumular los puntos de cada celda y promediarlos en tiempo $O(n)$.

\paragraph{Etapa 3: Tolerancia adaptativa.}
Si el parámetro de configuración \texttt{ground\_inlier\_tol} vale $0$ (valor
por defecto), se fija $\varepsilon$ automáticamente según la
ecuación~\ref{eq:adaptive-tol} a partir de $\sigma_{NN}$ calculado en la
Etapa 1. En caso contrario, el valor configurado se utiliza directamente,
lo que permite ajustar el comportamiento del algoritmo durante el desarrollo
sin modificar el código.

\paragraph{Etapa 4: RANSAC jerárquico.}
Se ejecuta el ciclo RANSAC (sección~\ref{subsec:lm-v3-ransac}) en dos fases.
La primera fase trabaja exclusivamente sobre los puntos candidatos a suelo,
definidos como aquellos cuya coordenada $Y \geq h_{\min}$, donde $h_{\min}$
es un parámetro configurable (por defecto $1{,}3\,\text{m}$).

El fundamento geométrico es el siguiente: en el frame alineado a gravedad,
$Y = 0$ corresponde a la posición de la cámara e $Y$ crece hacia abajo. La
cámara se monta en la cabeza del usuario, de modo que la distancia vertical
al suelo es aproximadamente la altura de la persona más un pequeño delta de
soporte. Como consecuencia, el suelo siempre se encuentra a $Y \geq h_{\min}$,
donde $h_{\min}$ es la altura mínima esperada de los usuarios del sistema.
Los puntos con $Y < h_{\min}$ están entre la cámara y el suelo (obstáculos,
cuerpo del usuario) y no constituyen candidatos válidos. El valor
$h_{\min} = 1{,}3\,\text{m}$ representa el peor caso conservador: si la
persona mide más, el filtro solo incorpora algunos puntos de obstáculos
adicionales en la fase 1, sin afectar la detección del plano.

Esta primera fase usa el $70\,$\% de las iteraciones disponibles. Si la
calidad del mejor candidato supera el umbral mínimo configurado (fracción de
inliers $\geq \rho_{\min}$, por defecto $0{,}3$), se transita directamente a
la Etapa 5. En caso contrario, la segunda fase amplía la búsqueda al conjunto
completo de puntos con las iteraciones restantes. Esta estrategia reduce el
número esperado de iteraciones en escenarios típicos con suelo visible.

\paragraph{Etapa 5: Validación geométrica.}
Se verifica que el plano candidato satisfaga dos criterios:
\begin{enumerate}
  \item El ángulo entre la normal del plano y el eje $Y$ es menor que
        $\alpha_{\max}$ (por defecto $15°$), condición que garantiza que
        el plano detectado es aproximadamente horizontal.
  \item La fracción de inliers supera $\rho_{\min}$.
\end{enumerate}
Un plano que no pase esta validación es descartado y la estimación devuelve
falla para ese fotograma.

\paragraph{Etapa 6: Refinamiento por mínimos cuadrados.}
Con el plano candidato como punto de partida, se colectan todos los puntos
de la nube \emph{completa} (no downsampled) cuya distancia al plano es
inferior a $1{,}5\,\varepsilon$ y se aplica el ajuste por mínimos cuadrados de la
sección~\ref{subsec:lm-v3-ls}. El plano refinado tiene menor sesgo porque
utiliza más puntos y no está restringido al subconjunto voxelizado.

\paragraph{Etapa 7: Clasificación.}
Cada punto de la nube original recibe una de las siguientes etiquetas,
determinadas exclusivamente por criterios geométricos:
\begin{itemize}
  \item \texttt{GROUND}: la distancia al plano refinado satisface
        $|\delta| < \varepsilon$. Estos puntos se publican en el tópico
        de suelo \texttt{/local\_mapper/debug/ground\_cloud}.

  \item \texttt{CEILING}: la coordenada satisface $Y_i <
        -\delta_{\text{techo}}$, donde $\delta_{\text{techo}}$ es un
        margen configurable (parámetro \texttt{ceiling\_delta\_m},
        valor por defecto $0{,}1\,\text{m}$). Dado que en el frame
        alineado $Y = 0$ corresponde a la posición de la cámara e $Y$
        crece hacia abajo, valores negativos de $Y$ indican puntos
        situados por encima de la cámara. Como la cámara se monta en
        la cabeza del usuario, estos puntos pertenecen al techo u
        otras superficies sobre la cabeza que no representan obstáculos
        chocables. Los puntos \texttt{CEILING} se descartan: no se
        añaden al tópico de obstáculos ni al mapa de ocupación.

  \item \texttt{OBSTACLE}: todo punto que no sea \texttt{GROUND} ni
        \texttt{CEILING}. Comprende paredes, objetos en el camino, el
        cuerpo del propio usuario y cualquier superficie entre la cámara
        y el suelo. Estos puntos se publican en
        \texttt{/local\_mapper/debug/obstacle\_cloud} y alimentan la
        grilla de ocupación.
\end{itemize}

\paragraph{Etapa 8: Análisis de performance.}
Se registran los tiempos de ejecución de cada etapa mediante
\texttt{clock\_gettime(CLOCK\_MONOTONIC)} para facilitar el ajuste del
presupuesto computacional en la plataforma destino (Raspberry Pi 5).

\subsection{Implementación}

\subsubsection{Interfaz de \texttt{GroundEstimator}}

La clase \texttt{GroundEstimator} expone la siguiente interfaz mínima:

\begin{itemize}
  \item \texttt{explicit GroundEstimator(const Config\& cfg)}: constructor
        que recibe una estructura de configuración con los parámetros de
        todas las etapas.
  \item \texttt{bool estimate(const std::vector<Point3D>\& cloud)}: ejecuta
        el pipeline completo y retorna \texttt{true} si se detectó un plano
        válido. Los resultados quedan disponibles en los getters.
  \item \texttt{const Plane\& groundPlane()}: parámetros del plano estimado
        (normal, distancia, calidad, número de inliers).
  \item \texttt{const std::vector<int>\& groundIndices()}: índices de los
        puntos clasificados como suelo.
  \item \texttt{const std::vector<int>\& obstacleIndices()}: índices de los
        puntos clasificados como obstáculo.
  \item \texttt{float cameraHeightM()}: distancia de la cámara al plano del
        suelo en metros, equivalente a $|d^*|$ de la
        ecuación~\ref{eq:plane-d}.
  \item \texttt{const CloudStats\& cloudStats()}: estadísticas de la Etapa 1.
  \item \texttt{const PerfStats\& perfStats()}: tiempos y conteos de la
        Etapa 8.
\end{itemize}

La estructura \texttt{Config} agrupa los parámetros de cada etapa y
proporciona valores por defecto conservadores para el caso de uso previsto:
cámara montada en la cabeza del usuario, por lo que la distancia vertical al
suelo es aproximadamente la altura de la persona más un pequeño delta de
soporte. El parámetro \texttt{min\_person\_height\_m} se fija en el peor caso
($1{,}3\,\text{m}$): si la persona resulta ser más alta, el filtro solo incluye
alguno puntos de obstáculos adicionales en la primera fase del RANSAC, sin
afectar la detección del plano del suelo.

\subsubsection{Integración en el nodo ROS}

La iteración V3 del nodo \texttt{DepthProjectionNode} añade los siguientes
elementos respecto a V2:

\begin{enumerate}
  \item Cinco nuevos parámetros ROS (\texttt{voxel\_size\_m},
        \texttt{ransac\_max\_iter}, \texttt{ground\_inlier\_tol},
        \texttt{min\_person\_height\_m}, \texttt{ceiling\_delta\_m}) que
        se leen de \texttt{params.yaml} en el constructor.
  \item Un objeto \texttt{GroundEstimator} inicializado con esos parámetros.
  \item El cuaternión de alineación gravitacional $q$ se calcula una única vez
        por mensaje IMU en \texttt{onImu()} y se almacena como miembro
        \texttt{q\_} de la clase. En \texttt{onDepth()}, ese valor cacheado
        se emplea para rotar explícitamente los puntos válidos al frame
        \texttt{gravity\_aligned\_frame}, produciendo la nube que se pasa
        a \texttt{ground\_estimator\_.estimate()}. Esta estrategia evita
        recalcular la orientación en cada frame de profundidad y mantiene
        coherencia temporal entre la nube y la estimación del plano.
  \item Cuatro nuevos publicadores de tipo \texttt{PointCloud2}:
        \texttt{ground\_cloud} (puntos clasificados como suelo),
        \texttt{obstacle\_cloud} (puntos clasificados como obstáculo),
        \texttt{ceiling\_cloud} (puntos descartados por estar sobre la cámara)
        y \texttt{raw\_cloud} (nube cruda en \texttt{gravity\_aligned\_frame},
        sin rotar, para verificación visual del alineamiento).
        Los primeros tres se publican en \texttt{gravity\_aligned\_frame}
        con coordenadas ya rotadas; \texttt{raw\_cloud} se publica con
        coordenadas de cámara directamente en ese frame, de modo que en
        RViz se observe la oscilación al inclinar el dispositivo.
        La nube cruda principal (\texttt{depth\_cloud}) se publica en
        \texttt{camera\_depth\_optical\_frame}: RViz aplica la TF dinámica
        publicada en \texttt{onImu()} para mostrarla alineada, sin necesidad
        de una copia rotada explícita.
\end{enumerate}

La rotación de los puntos al frame alineado se aplica mediante el método
\texttt{Quaternion::rotatePoint()} de \texttt{GravityAligner}
(véase la sección~\ref{sec:lm-gravity-aligner}), que implementa la
operación $\mathbf{p}' = q\,\mathbf{p}\,q^{*}$ con aritmética de
cuaterniones unitarios en punto flotante de simple precisión.

\subsubsection{Ausencia de dependencias externas}

La implementación de \texttt{GroundEstimator} es C++17 puro, sin dependencias
de ROS ni de bibliotecas de álgebra lineal (PCL, Eigen). El algoritmo de
Jacobi para el cálculo de autovectores $3 \times 3$, la tabla hash para
el voxel grid y el generador de números pseudoaleatorios de la biblioteca
estándar son suficientes para el problema. Esto facilita el testing unitario
mediante GTest sin necesitar el entorno de ROS.

\subsection{Experimentos propuestos}

Los siguientes experimentos permiten evaluar el desempeño del sistema en
condiciones reales antes de integrar la detección del suelo con el mapeador
de ocupación.

\paragraph{Superficie plana perfecta.}
Apuntar la cámara hacia un suelo liso (azulejo o tablero) a distintas
distancias ($0{,}5\,\text{m}$, $1\,\text{m}$, $2\,\text{m}$). Verificar
que el pipeline estima un plano con calidad superior a $0{,}6$ y que la
altura estimada $h$ coincide con la medición manual en $\pm 5\,\text{cm}$.

\paragraph{Inclinación deliberada.}
Inclinar la cámara $\pm 20°$ respecto de la horizontal sin mover el
dispositivo del mismo punto. Con la alineación gravitacional activa (V2),
el plano del suelo debe seguir siendo detectado correctamente. Sin ella,
la Etapa 5 debería rechazar el plano.

\paragraph{Escena con obstáculos.}
Colocar objetos (caja, silla) frente a la cámara y verificar que los puntos
de tales objetos no son clasificados como \texttt{GROUND}.

\paragraph{Benchmark de performance.}
Medir \texttt{PerfStats.time\_total\_ms} durante $100$ fotogramas con la
cámara a $30\,\text{fps}$. El presupuesto objetivo es $< 15\,\text{ms}$ por
fotograma en Raspberry Pi 5, lo que deja margen para el resto del pipeline
de navegación.
